\documentclass[12pt]{article}
\usepackage[margin=1in]{geometry}
\usepackage{amsmath,amssymb,amsthm}
\usepackage{booktabs}
\usepackage[colorlinks=true,linkcolor=blue,citecolor=blue,urlcolor=blue]{hyperref}

\newcommand{\seqnum}[1]{\href{https://oeis.org/#1}{\underline{#1}}}
\newcommand{\fl}[1]{\lfloor #1 \rfloor}
\newcommand{\cl}[1]{\lceil #1 \rceil}
\newcommand{\Z}{\mathbb{Z}}
\newcommand{\R}{\mathbb{R}}
\newcommand{\D}{\mathcal{D}}
\newcommand{\M}{\mathcal{M}}
\newcommand{\Ws}{W^{*}}

\theoremstyle{plain}
\newtheorem{theorem}{Theorem}
\newtheorem{lemma}[theorem]{Lemma}
\newtheorem{proposition}[theorem]{Proposition}
\newtheorem{corollary}[theorem]{Corollary}
\theoremstyle{remark}
\newtheorem{remark}[theorem]{Remark}

\title{Numbers Occurring Twice in, and Missing from,\\ the Wythoff Ceiling Array}
\author{Camden Link}
\date{Draft of November 21, 2026}

\begin{document}
\maketitle

\begin{abstract}
Kimberling's Wythoff ceiling array is obtained from a formula for the Wythoff array by
replacing a floor with a ceiling. Unlike the Wythoff array, it contains some positive integers
twice and misses others. In 2021 Kimberling conjectured that consecutive numbers occurring twice
differ by $3$, $5$, $10$ or $13$, and that consecutive missing numbers differ by $2$, $3$, $5$,
$10$ or $13$. We prove both conjectures, together with the statement, recorded without proof in
the On-Line Encyclopedia of Integer Sequences, that no number occurs more than twice. The gap $2$
between missing numbers occurs only once, and every other gap occurs infinitely often. We also
determine the densities and gap frequencies of both sets, and describe both sets by Zeckendorf
representations. The proof reduces membership in a column of the array to the position of a
number modulo the golden ratio. The main results are formally verified in the Lean proof
assistant.
\end{abstract}

\section{Introduction}

Let $\varphi=(1+\sqrt5)/2$ and let $F(0)=0$, $F(1)=1$, $F(k+1)=F(k)+F(k-1)$ be the Fibonacci
numbers. The Wythoff array \seqnum{A035513} is
\[
W(n,k)=F(k+1)\fl{n\varphi}+(n-1)F(k)\qquad(n,k\ge1).
\]
Each of its rows satisfies the Fibonacci recurrence, and it contains every positive integer
exactly once \cite{morrison,kimberling95}. In February 2021 Kimberling \cite{oeis} introduced
the \emph{Wythoff ceiling array} \seqnum{A341260},
\[
\Ws(n,k)=F(k+1)\cl{n\varphi}+(n-1)F(k)\qquad(n,k\ge1),
\]
whose corner is
\[
\begin{array}{rrrrrrr}
2&4&6&10&16&26&\cdots\\
5&9&14&23&37&60&\cdots\\
7&12&19&31&50&81&\cdots\\
10&17&27&44&71&115&\cdots
\end{array}
\]
Each row still satisfies $\Ws(n,k+1)=\Ws(n,k)+\Ws(n,k-1)$, but the array no longer contains
each positive integer once: $10=\Ws(1,4)=\Ws(4,1)$ occurs twice, while $1$, $3$ and $8$ do not
occur. Kimberling recorded the two sets
\begin{align*}
\D&=\{10,23,26,31,44,57,60,65,68,78,\dots\}&&\text{numbers occurring twice (\seqnum{A341261}),}\\
\M&=\{1,3,8,11,21,24,29,32,45,55,\dots\}&&\text{positive integers not occurring (\seqnum{A341262}).}
\end{align*}
The entries \seqnum{A341260} and \seqnum{A341261} state that no number occurs more than twice,
without proof. They also record two conjectures:
\begin{quote}
(C1) consecutive terms of $\D$ differ only by $3$, $5$, $10$ or $13$;\\
(C2) consecutive terms of $\M$ differ only by $2$, $3$, $5$, $10$ or $13$.
\end{quote}
As of November 2026 the entries still list both as conjectures, and we have not found them
settled elsewhere. We call $g$ a \emph{gap} of a set $S\subseteq\Z$ if $m$ and $m+g$ are
consecutive elements of $S$ for some $m$.

\begin{theorem}\label{thm:main}
\leavevmode
\begin{enumerate}
\item[(a)] No positive integer occurs at more than two positions $(n,k)$ of the array.
\item[(b)] The gaps of $\D$ are exactly $3,5,10,13$, and each of them occurs infinitely often.
\item[(c)] The gaps of $\M$ are exactly $2,3,5,10,13$. The gap $2$ occurs only between $1$ and
$3$; each of the others occurs infinitely often.
\end{enumerate}
\end{theorem}

In particular both conjectures hold, and every value they allow does occur.

\begin{theorem}\label{thm:freq}
The sets $\D$ and $\M$ each have density $\varphi^{-4}=(7-3\sqrt5)/2\approx0.1459$. In each set
the gaps $3$, $5$, $10$, $13$ occur with relative frequencies $\varphi^{-2}$, $\varphi^{-3}$,
$\varphi^{-4}$, $\varphi^{-3}$, that is, about $38.2\%$, $23.6\%$, $14.6\%$, $23.6\%$.
\end{theorem}

Every positive integer $m$ has a unique \emph{Zeckendorf representation}
$m=F(i_1)+F(i_2)+\cdots+F(i_r)$ with $2\le i_1$ and $i_{s+1}\ge i_s+2$ \cite{zeckendorf}.

\begin{theorem}\label{thm:zeck}
Let $m\ge1$ have Zeckendorf representation $F(i_1)+\cdots+F(i_r)$ with $i_1<\cdots<i_r$.
\begin{enumerate}
\item[(a)] $m\in\D$ if and only if $r\ge2$, $i_1$ is odd and $i_2$ is even.
\item[(b)] $m\in\M$ if and only if $m=1$, or $i_1\ge4$ is even and either $r=1$ or $i_2$ is even.
\end{enumerate}
\end{theorem}

For example $10=F(3)+F(6)\in\D$ and $11=F(4)+F(6)\in\M$. In words: split the Fibonacci numbers
into $2,5,13,34,\dots$ (odd index) and $1,3,8,21,\dots$ (even index); a number occurs twice when
its smallest Zeckendorf part is of the first kind and its second smallest is of the second kind.

\subsection*{Method}

The key observation (Proposition~\ref{prop:col}) is that $m$ lies in column $k$ exactly when
$m-\varphi x$ lies in the interval $\psi^k\cdot(\varphi,\varphi+1)$ for some integer $x$, where
$\psi=1-\varphi=-1/\varphi$. These intervals are disjoint and shrink geometrically towards $0$,
alternating sides. So whether $m$ occurs, and how often, depends only on $m$ modulo $\varphi$,
and passing from $m$ to $m+1$ rotates the circle $\R/\varphi\Z$ by $1$. The gaps are then
return times of this rotation to a union of infinitely many intervals accumulating at one point.
Only the first three of these intervals have to be treated separately; the rest are handled
together. The approach, reducing membership to the position of a point on a circle and checking
finitely many cells by machine, is the one used in \cite{link} for intersections of Beatty
sequences.

\subsection*{Formal verification}

Theorem~\ref{thm:main} is formally verified in Lean~4 \cite{lean4} with Mathlib \cite{mathlib},
as are Proposition~\ref{prop:col}, Theorem~\ref{thm:struct}(a) and the parts of
Theorem~\ref{thm:struct}(b),(c) used in Section~\ref{sec:gaps}. Theorems~\ref{thm:freq}
and~\ref{thm:zeck} are proved here in the usual way and are not formalized; both were also
checked against a direct enumeration of the array up to $2\cdot10^9$.
Section~\ref{sec:lean} says precisely what is machine-checked. The Lean sources and the programs
are available at \url{https://github.com/clink-bw/math-proofs}.

\subsection*{Organization}

Section~\ref{sec:prelim} collects facts about $\varphi$. Section~\ref{sec:circle} describes
the columns on the circle and proves Theorems~\ref{thm:main}(a) and~\ref{thm:zeck}.
Section~\ref{sec:gaps} proves Theorem~\ref{thm:main}(b),(c) and Theorem~\ref{thm:freq}.
Section~\ref{sec:lean} describes the formalization, Section~\ref{sec:remarks} makes some
remarks, and Section~\ref{sec:ai} states how AI tools were used.

\section{Preliminaries}\label{sec:prelim}

Throughout, $\psi=1-\varphi=-1/\varphi$, so that $\varphi\psi=-1$, $\varphi^2=\varphi+1$ and
$|\psi|=\varphi^{-1}<1$. Binet's formula gives
\begin{equation}\label{eq:binet}
F(k+1)-\varphi F(k)=\psi^k\qquad(k\ge0).
\end{equation}
Hence each $\varphi^{-j}=|\psi|^j=(-1)^j\bigl(F(j+1)-\varphi F(j)\bigr)$ lies in $\Z+\Z\varphi$:
\[
\varphi^{-1}=\varphi-1,\quad \varphi^{-2}=2-\varphi,\quad \varphi^{-3}=2\varphi-3,\quad
\varphi^{-4}=5-3\varphi,\quad \varphi^{-5}=5\varphi-8,\quad \varphi^{-6}=13-8\varphi,
\]
and $\varphi^{-j}=\varphi^{-j-1}+\varphi^{-j-2}$. Since $\varphi$ is irrational,
$a+b\varphi=a'+b'\varphi$ with integers $a,b,a',b'$ forces $a=a'$ and $b=b'$.

For $k\ge0$ let $I_k=\{\psi^k u:\ \varphi<u<\varphi+1\}$. As $\psi^k\varphi=-\psi^{k-1}$ and
$\psi^k(\varphi+1)=\psi^k\varphi^2=\psi^{k-2}$,
\begin{equation}\label{eq:Ik}
I_k=(\varphi^{1-k},\ \varphi^{2-k})\ \text{ for $k$ even},\qquad
I_k=(-\varphi^{2-k},\ -\varphi^{1-k})\ \text{ for $k$ odd}.
\end{equation}
Thus $I_1=(-\varphi,-1)$, $I_2=(\varphi^{-1},1)$, $I_3=(-\varphi^{-1},-\varphi^{-2})$,
$I_4=(\varphi^{-3},\varphi^{-2})$, and so on. The intervals $I_k$ are pairwise disjoint; for even
$k\ge2$ they lie in $(0,1)$ and for odd $k$ in $(-\varphi,0)$.

For an integer $m$ let
\[
t(m)=m-\varphi\fl{m/\varphi}\in[0,\varphi),
\]
the position of $m$ on the circle $\R/\varphi\Z$. Then $t(m+j)\equiv t(m)+j\pmod\varphi$, and
\begin{equation}\label{eq:irr}
\text{if } t(m)=a+b\varphi \text{ with } a,b\in\Z, \text{ then } a=m.
\end{equation}
In particular $t(m)\neq0$ for $m\ne0$.

\section{The columns on the circle}\label{sec:circle}

It is convenient to extend the array by $\Ws(n,0)=\cl{n\varphi}$, which is the defining formula
with $k=0$.

\begin{lemma}\label{lem:row}
For $n\ge1$ and $k\ge0$,
\[
\Ws(n,k+1)-\varphi\,\Ws(n,k)=\psi^{k+1}u_n,\qquad u_n=\cl{n\varphi}-(n-1)\varphi,
\]
and $\varphi<u_n<\varphi+1$.
\end{lemma}

\begin{proof}
By \eqref{eq:binet} the left side is
$\cl{n\varphi}\bigl(F(k+2)-\varphi F(k+1)\bigr)+(n-1)\bigl(F(k+1)-\varphi F(k)\bigr)
=\cl{n\varphi}\psi^{k+1}+(n-1)\psi^k$, and $(n-1)\psi^k=-(n-1)\varphi\psi^{k+1}$. Since
$n\varphi$ is irrational, $n\varphi<\cl{n\varphi}<n\varphi+1$, which gives the bounds on $u_n$.
\end{proof}

\begin{proposition}\label{prop:col}
Let $m\ge1$, $k\ge1$ and $x\in\Z$. Then $m-\varphi x\in I_k$ if and only if there is $n\ge1$ with
$\Ws(n,k)=m$ and $\Ws(n,k-1)=x$. In particular $m$ lies in column $k$ if and only if
$m-\varphi x\in I_k$ for some integer $x$.
\end{proposition}

\begin{proof}
One direction is Lemma~\ref{lem:row}. For the other we use induction on $k$. Let $k=1$, so
$-\varphi<m-\varphi x<-1$, and put $n=m-x+1$. Multiplying by $\varphi$ and using
$\varphi^2=\varphi+1$ gives $x-1<n\varphi<x$, so $\cl{n\varphi}=x=\Ws(n,0)$ and
$\Ws(n,1)=\cl{n\varphi}+n-1=m$. Also $n\ge1$: if $x\le0$ this is clear, and if $x\ge1$ then
$m>\varphi(x-1)\ge x-1$. Now let $k\ge2$. By \eqref{eq:Ik}, $|m-\varphi x|<1$, so
$\varphi x>m-1\ge0$ and $x\ge1$. Moreover
$x-\varphi(m-x)=-\varphi(m-\varphi x)\in-\varphi I_k=I_{k-1}$. By induction there is $n\ge1$ with
$\Ws(n,k-1)=x$ and $\Ws(n,k-2)=m-x$, and then $\Ws(n,k)=\Ws(n,k-1)+\Ws(n,k-2)=m$.
\end{proof}

Since $\cl{(n+1)\varphi}\ge\cl{n\varphi}+1$, each column is strictly increasing in $n$. So a
number occurs at most once in each column, and the number of times it occurs is the number of
columns containing it.

\begin{corollary}\label{cor:t}
Let $m\ge1$ and $t=t(m)$. For $k\ge1$, $m$ lies in column $k$ if and only if $t\in I_k$ (for
even $k$) or $t-\varphi\in I_k$ (for odd $k$).
\end{corollary}

\begin{proof}
Each $m-\varphi x$ equals $t+\varphi(\fl{m/\varphi}-x)$, and all $I_k$ lie in $(-\varphi,1)$.
As $0\le t<\varphi$, the only values in $(-\varphi,1)$ are $t\ge0$ and $t-\varphi<0$; the first
can lie only in an even $I_k$ and the second only in an odd one.
\end{proof}

\begin{theorem}\label{thm:struct}
Let $m\ge1$ and $t=t(m)$, and put
\[
T=\bigcup_{j\ge1}\bigl(\varphi^{-2j-1},\ \varphi^{-2j}\bigr),\qquad
T'=\bigcup_{j\ge1}\bigl[\varphi^{-2j-1},\ \varphi^{-2j}\bigr).
\]
\begin{enumerate}
\item[(a)] $m$ occurs at most twice in the array.
\item[(b)] $m\in\D$ if and only if $m$ lies in column $1$ and in some even column $k\ge4$, if and
only if $t\in T$.
\item[(c)] $m\in\M$ if and only if $t=1$ (that is, $m=1$) or $\varphi-t\in T'$.
\end{enumerate}
\end{theorem}

\begin{proof}
By Corollary~\ref{cor:t} the number of occurrences of $m$ is
$[t\in E]+[t\in O]$, where
\begin{align*}
E&=\bigcup_{k\ge2\text{ even}}I_k=\bigcup_{j\ge1}(\varphi^{1-2j},\ \varphi^{2-2j})\subset(0,1),\\
O&=\bigcup_{k\text{ odd}}(I_k+\varphi)=(0,\varphi^{-1})\cup
\bigcup_{j\ge1}(\varphi-\varphi^{1-2j},\ \varphi-\varphi^{-2j}).
\end{align*}
The last union lies in $(1,\varphi)$. Since the $I_k$ are disjoint, each of $E$ and $O$
contributes at most one column, which proves (a). For (b), $E\cap O=E\cap(0,\varphi^{-1})$,
which is $T$; there $m$ lies in column $1$ (from $I_1+\varphi=(0,\varphi^{-1})$) and in column
$2j+2$ (from $I_{2j+2}$). For (c), $E\cup O$ contains $(0,1)$ except $\varphi^{-1}$. By
\eqref{eq:irr}, $t=\varphi^{-1}=\varphi-1$ would give $m=-1$; hence $t\ge1$, and $t=1$ gives
$m=1$, which is indeed missing. For $1<t<\varphi$ we have $t\notin E$, and $t\notin O$ if and
only if $\varphi-t$, which lies in $(0,\varphi^{-1})$, avoids every
$(\varphi^{-2j},\varphi^{1-2j})$, that is, $\varphi-t\in\bigcup_{j\ge1}[\varphi^{-2j-1},\varphi^{-2j}]$.
The right endpoints do not occur: $\varphi-t=\varphi^{-2j}$ would give
$t=\varphi-F(2j+1)+F(2j)\varphi$ and $m=-F(2j+1)$ by \eqref{eq:irr}.
\end{proof}

\begin{remark}
The left endpoints of $T'$ do occur. By \eqref{eq:binet},
$F(i)=\varphi F(i-1)+\psi^{i-1}$, so $t(F(2j+2))=\varphi-\varphi^{-2j-1}$ and the even-indexed
Fibonacci numbers $3,8,21,55,\dots$ are missing. Likewise $t(F(2j+1))=\varphi^{-2j}$ is an
endpoint of $T$, and the odd-indexed Fibonacci numbers $2,5,13,\dots$ occur exactly once.
\end{remark}

\begin{proof}[Proof of Theorem~\ref{thm:zeck}]
Call a finite set $J$ of nonnegative integers \emph{sparse} if it contains no two consecutive
integers. For sparse $J$, the sum $\Sigma(J)=\sum_{j\in J}\psi^j$ lies strictly between the sum
over all odd $j$, which is $\psi/(1-\psi^2)=-1$, and the sum over all even $j$, which is
$1/(1-\psi^2)=\varphi$; here $1-\psi^2=-\psi$. If moreover $0\in J$, then
$\Sigma(J)=1+\psi^2\Sigma(J')$ with $J'$ sparse, so $\Sigma(J)\in(\varphi^{-1},\varphi)$.

Let $m=F(i_1)+\cdots+F(i_r)$. By \eqref{eq:binet},
\[
m-\varphi\sum_s F(i_s-1)=\sum_s\psi^{i_s-1}=:\sigma=\psi^{i_1-1}\tau,\qquad
\tau=\sum_s\psi^{i_s-i_1}\in(\varphi^{-1},\varphi).
\]
Since $\sigma=\psi\,\Sigma(\{i_s-2\})\in\psi\cdot(-1,\varphi)=(-1,\varphi^{-1})$, we have
$t(m)=\sigma$ if $\sigma>0$ and $t(m)=\sigma+\varphi$ if $\sigma<0$. If $r=1$ then $\tau=1$.
If $r\ge2$, let $d=i_2-i_1\ge2$; then $\tau-1=\psi^d\rho$ with $\rho\in(\varphi^{-1},\varphi)$,
so $\tau<1$ if $d$ is odd and $\tau>1$ if $d$ is even.

(a) As $T\subset(0,\varphi^{-2})$ and $\sigma+\varphi>\varphi^{-1}$, membership in $\D$ requires
$\sigma>0$, that is, $i_1=2a+1$ odd, with $a\ge1$. Then $t=\varphi^{-2a}\tau$ lies in
$(\varphi^{-2a-1},\varphi^{1-2a})$, which meets $T$ only in $(\varphi^{-2a-1},\varphi^{-2a})$. So
$m\in\D$ if and only if $\tau<1$, that is, $r\ge2$ and $d$ odd, that is, $i_2$ even.

(b) If $\sigma>0$ then $t<\varphi^{-1}$ and $m\notin\M$ by Theorem~\ref{thm:struct}(c). So let
$i_1=2a$ be even; then $\varphi-t=-\sigma=\varphi^{1-2a}\tau$. If $a=1$ this lies in
$(\varphi^{-2},1)$, which misses $T'$, and it equals $\varphi^{-1}$ (that is, $t=1$) exactly when
$\tau=1$, that is, $m=F(2)=1$. If $a\ge2$, then $t\ne1$ and $\varphi-t\in(\varphi^{-2a},\varphi^{2-2a})$,
which meets $T'$ only in $[\varphi^{1-2a},\varphi^{2-2a})$. So $m\in\M$ if and only if
$\tau\ge1$, that is, $r=1$ or $d$ even, that is, $r=1$ or $i_2$ even.
\end{proof}

\section{The gaps}\label{sec:gaps}

By Theorem~\ref{thm:struct}, a gap of $\D$ is a return time of the rotation
$t\mapsto t+1\pmod\varphi$ to $T$, and similarly for $\M$. Table~\ref{tab:rot} lists
$\rho_j=j-\varphi\fl{j/\varphi}$, so that $t(m+j)$ is $t(m)+\rho_j$ or $t(m)+\rho_j-\varphi$.

\begin{table}[ht]
\centering
\small
\setlength{\tabcolsep}{4pt}
\begin{tabular}{clc@{\qquad}clc}
\toprule
$j$ & $\rho_j$ & $\approx$ & $j$ & $\rho_j$ & $\approx$ \\
\midrule
$1$ & $1$ & 1.000 & $8$ & $8-4\varphi$ & 1.528 \\
$2$ & $2-\varphi$ & 0.382 & $9$ & $9-5\varphi$ & 0.910 \\
$3$ & $3-\varphi$ & 1.382 & $10$ & $10-6\varphi$ & 0.292 \\
$4$ & $4-2\varphi$ & 0.764 & $11$ & $11-6\varphi$ & 1.292 \\
$5$ & $5-3\varphi$ & 0.146 & $12$ & $12-7\varphi$ & 0.674 \\
$6$ & $6-3\varphi$ & 1.146 & $13$ & $13-8\varphi$ & 0.056 \\
$7$ & $7-4\varphi$ & 0.528 & & & \\
\bottomrule
\end{tabular}
\caption{The rotation amounts $\rho_j=j\bmod\varphi$. Note $\rho_2=\varphi^{-2}$,
$\rho_3=\varphi-\varphi^{-3}$, $\rho_5=\varphi^{-4}$, $\rho_8=\varphi-\varphi^{-5}$,
$\rho_{10}=2\varphi^{-4}$, $\rho_{13}=\varphi^{-6}$.}
\label{tab:rot}
\end{table}

\begin{lemma}\label{lem:shift}
For $m\in\Z$ and $i\ge0$, $t(m+F(i+1))\equiv t(m)+\psi^i\pmod\varphi$.
\end{lemma}

\begin{proof}
By \eqref{eq:binet}, $(m+F(i+1))-\varphi(\fl{m/\varphi}+F(i))=t(m)+\psi^i$.
\end{proof}

Since $\psi^i\to0$ and $F(i+1)\to\infty$, if $t(m)$ lies in an open interval $C\subset(0,\varphi)$
then $t(m')\in C$ for infinitely many $m'>m$.

% Generated by wythoff-ceiling/compute/paper_tables.py; do not edit.
\begin{table}[ht]
\centering
\small
\setlength{\tabcolsep}{3.2pt}
\begin{tabular}{llccccccccccccc}
\toprule
cell of $t(m)$ & gap & $1$ & $2$ & $3$ & $4$ & $5$ & $6$ & $7$ & $8$ & $9$ & $10$ & $11$ & $12$ & $13$ \\
\midrule
$(\varphi^{-3}+\varphi^{-5},\ \varphi^{-2})$ & 3 & a & a & $\mathbf{6}$ &  &  &  &  &  &  &  &  &  &  \\
$(\varphi^{-3}+\varphi^{-6},\ \varphi^{-3}+\varphi^{-5})$ & 13 & a & a & c & a & a & a & a & b & a & a & a & a & $\mathbf{4}$ \\
$(\varphi^{-3},\ \varphi^{-3}+\varphi^{-6}]$ & 3 or 13 & a & a & ? & a & a & a & a & b & a & a & a & a & $\mathbf{4}$ \\
$(\varphi^{-5},\ \varphi^{-4})$ & 5 & a & a & a & a & $\mathbf{4}$ &  &  &  &  &  &  &  &  \\
$(0,\ \varphi^{-6})$ & 10 & a & a & a & a & b & a & a & a & a & $\mathbf{4}$ &  &  &  \\
\bottomrule
\end{tabular}
\caption{Why $m+j\in\mathcal D$ or $m+j\notin\mathcal D$, for $m$ with $t(m)$ in the given cell. A bold $k$: $t(m+j)\in I_k$, so $m+j$ is in columns $1$ and $k$. a: $t(m+j)\ge\varphi^{-2}$. b: $t(m+j)\in[\varphi^{-4},\varphi^{-3}]$. c: $t(m+j)\in[\varphi^{-6},\varphi^{-5}]$. In cases a, b, c no even column $\ge4$ contains $m+j$. ?: either, see the text.}
\label{tab:twice}
\end{table}

\begin{table}[ht]
\centering
\small
\setlength{\tabcolsep}{3.2pt}
\begin{tabular}{llccccccccccccc}
\toprule
cell of $\varphi-t(m)$ & gap & $1$ & $2$ & $3$ & $4$ & $5$ & $6$ & $7$ & $8$ & $9$ & $10$ & $11$ & $12$ & $13$ \\
\midrule
$\{\varphi^{-1}\}$ & 2 & 1 & $\ast$ &  &  &  &  &  &  &  &  &  &  &  \\
$[\varphi^{-3},\ 2\varphi^{-4})$ & 5 & 2 & 1 & 3 & 1 & $\ast$ &  &  &  &  &  &  &  &  \\
$(2\varphi^{-4},\ \varphi^{-2})$ & 10 or 13 & 2 & 1 & 3 & 1 & 5 & 2 & 1 & 3 & 1 & ? & 2 & 1 & $\ast$ \\
$[\varphi^{-5},\ \varphi^{-4})$ & 3 & 2 & 1 & $\ast$ &  &  &  &  &  &  &  &  &  &  \\
$(0,\ \varphi^{-6})$ & 3 & 2 & 1 & $\ast$ &  &  &  &  &  &  &  &  &  &  \\
\bottomrule
\end{tabular}
\caption{Why $m+j\in\mathcal M$ or $m+j\notin\mathcal M$, for $m$ with $\varphi-t(m)$ in the given cell. A number $k$: $m+j$ is in column $k$. $\ast$: $m+j\in\mathcal M$. ?: either, see the text.}
\label{tab:missing}
\end{table}


\begin{proof}[Proof of Theorem~\ref{thm:main}(b)]
Let $m\in\D$ and $t=t(m)\in T$. Since $T\cap[\varphi^{-6},\varphi)=I_4\cup I_6$, and
$t\ne\varphi^{-3}+\varphi^{-5}=7\varphi-11$ by \eqref{eq:irr}, $t$ lies in one of the five cells
of Table~\ref{tab:twice}. For each cell and each $j$, the table records where $t(m+j)$ lies,
computed from Table~\ref{tab:rot}. By Theorem~\ref{thm:struct}(b), $m+j\in\D$ exactly when
$t(m+j)\in T$. The entries a, b, c place $t(m+j)$ outside $T$, which lies in $(0,\varphi^{-2})$
and avoids $[\varphi^{-4},\varphi^{-3}]$ and $[\varphi^{-6},\varphi^{-5}]$. A bold $4$ or $6$
places $t(m+j)$ in $I_4$ or $I_6$, which lie in $T$. For example, if $t$ is in the first cell then
$t(m+3)=t-\varphi^{-3}\in(\varphi^{-5},\varphi^{-4})=I_6$.

In the third cell, $t(m+3)=t-\varphi^{-3}$ lies in $(0,\varphi^{-6}]$, which meets both $T$ and its
complement. If $m+3\in\D$ the gap is $3$. Otherwise the remaining entries show that $m+13$ is the
next element: $t(m+13)=t+\varphi^{-6}\in(\varphi^{-3}+\varphi^{-6},\varphi^{-3}+2\varphi^{-6}]\subset I_4$.
Hence the next element of $\D$ after $m$ is $m+g$ with $g\in\{3,5,10,13\}$.

In the first, second and fourth cells, and in $I_8=(\varphi^{-7},\varphi^{-6})\subset T$ inside the
fifth, the gap is $3$, $13$, $5$ and $10$ respectively at every point. These cells contain $t(23)$,
$t(44)$, $t(26)$ and $t(68)$, which are approximately $0.348$, $0.313$, $0.111$ and $0.043$. By
Lemma~\ref{lem:shift} each gap occurs infinitely often.
\end{proof}

\begin{proof}[Proof of Theorem~\ref{thm:main}(c)]
Let $m\in\M$ and $v=\varphi-t(m)$. By Theorem~\ref{thm:struct}(c), either $m=1$ and
$v=\varphi^{-1}$, or $v\in T'$. In the second case $v$ lies in one of the last four cells of
Table~\ref{tab:missing}, since $T'\cap[\varphi^{-6},\varphi)=[\varphi^{-3},\varphi^{-2})\cup
[\varphi^{-5},\varphi^{-4})$, $2\varphi^{-4}=\varphi^{-3}+\varphi^{-6}$, and
$v\ne2\varphi^{-4}$ because $t(m)=7\varphi-10$ would give $m=-10$. An entry $k$ in the table means
that $m+j$ lies in column $k$, by Corollary~\ref{cor:t}, so $m+j\notin\M$. An entry $\ast$ means
that $\varphi-t(m+j)$ lies in $[\varphi^{-3},\varphi^{-2})$, $[\varphi^{-5},\varphi^{-4})$ or
$[\varphi^{-7},\varphi^{-6})$, all contained in $T'$, so $m+j\in\M$.

In the third row, $\varphi-t(m+10)=v-2\varphi^{-4}\in(0,\varphi^{-5})$. If $m+10\in\M$ the gap is
$10$; otherwise it is $13$, as $\varphi-t(m+13)=v-\varphi^{-6}\in(\varphi^{-3},\varphi^{-2})$.
The gap $2$ arises only in the first row, that is, for $m=1$. The gaps $5$, $3$, $10$ and $13$
occur at every point of $(\varphi^{-3},2\varphi^{-4})$, of $(\varphi^{-5},\varphi^{-4})$, of
$2\varphi^{-4}+(\varphi^{-7},\varphi^{-6})$ and of $2\varphi^{-4}+(\varphi^{-6},\varphi^{-5})$
respectively; these contain $\varphi-t(m)$ for $m=24$, $29$, $66$ and $32$. By
Lemma~\ref{lem:shift} each of them occurs infinitely often.
\end{proof}

\begin{proof}[Proof of Theorem~\ref{thm:freq}]
The sequence $t(m)=\varphi\{m/\varphi\}$ is equidistributed in $[0,\varphi)$ by Weyl's theorem
\cite{weyl}. The sets $T$, $T'$ and the cells above are unions of intervals whose boundary is a
countable set with one accumulation point, so their indicator functions are Riemann integrable,
and the proportion of $m\le N$ with $t(m)$ in such a set tends to its length divided by $\varphi$
\cite{kn}. Since $\varphi^{-2j}-\varphi^{-2j-1}=\varphi^{-2j-2}$,
\[
|T|=|T'|=\sum_{j\ge1}\varphi^{-2j-2}=\frac{\varphi^{-4}}{1-\varphi^{-2}}=\varphi^{-3},
\]
which gives density $\varphi^{-4}$ for both sets. For $\D$, the gap is $5$ on $I_6$ (length
$\varphi^{-6}$); $10$ on $T\cap(0,\varphi^{-6})$ (length $\sum_{j\ge3}\varphi^{-2j-2}=\varphi^{-7}$);
$3$ on $\{t\in I_4:\ t-\varphi^{-3}\in T\}$ (length $|T\cap(0,\varphi^{-4})|=\varphi^{-5}$); and
$13$ on the rest of $I_4$ (length $\varphi^{-4}-\varphi^{-5}=\varphi^{-6}$). Dividing by
$|T|=\varphi^{-3}$ gives the frequencies $\varphi^{-3},\varphi^{-4},\varphi^{-2},\varphi^{-3}$.
For $\M$, the gap is $5$ on $[\varphi^{-3},2\varphi^{-4})$ (length $\varphi^{-6}$); $3$ on
$[\varphi^{-5},\varphi^{-4})$ and on $T'\cap(0,\varphi^{-6})$ (total length
$\varphi^{-6}+\varphi^{-7}=\varphi^{-5}$); $10$ on $\{v\in(2\varphi^{-4},\varphi^{-2}):\ v-2\varphi^{-4}\in T'\}$
(length $\varphi^{-7}$); and $13$ on the rest of $(2\varphi^{-4},\varphi^{-2})$ (length
$\varphi^{-5}-\varphi^{-7}=\varphi^{-6}$). The frequencies are the same.
\end{proof}

A direct enumeration up to $2\cdot10^9$ gives $291{,}796{,}059$ elements of $\D$ and
$291{,}796{,}079$ of $\M$, and relative gap frequencies agreeing with Theorem~\ref{thm:freq} to
six decimal places.

\section{Formal verification}\label{sec:lean}

\begin{sloppypar}
The following statements are proved in Lean~4 with Mathlib, in the namespace
\texttt{WythoffCeiling}. The development defines the array as a function \texttt{W} from pairs of
natural numbers to integers by the formula above, $\D$ as the set of integers
$m$ with $\texttt{W}\,n\,k=m=\texttt{W}\,n'\,k'$ for two different pairs with
$n,k,n',k'\ge1$, and $\M$ as the set of positive integers not of the form $\texttt{W}\,n\,k$
with $n,k\ge1$. As in \cite{link}, \texttt{IsGap S g} means that some $m\in S$ has $m+g\in S$ and
$m+j\notin S$ for $0<j<g$.
\end{sloppypar}

\begin{itemize}
\item \texttt{at\_most\_twice}: Theorem~\ref{thm:main}(a).
\item \texttt{isGap\_twice\_iff}: for $g>0$, \texttt{IsGap} $\D$ $g$ holds if and only if
$g\in\{3,5,10,13\}$; \texttt{twice\_gap\_infinitely\_often}: for each such $g$ and each bound $B$
there is $m>B$ in $\D$ whose successor in $\D$ is $m+g$. Together, Theorem~\ref{thm:main}(b).
\item \texttt{isGap\_missing\_iff}, \texttt{missing\_gap\_infinitely\_often} and
\texttt{missing\_two\_apart} (if $m$ and $m+2$ are both in $\M$ then $m=1$): together,
Theorem~\ref{thm:main}(c).
\item \texttt{inCol\_iff}: Proposition~\ref{prop:col}; \texttt{colR\_iff}:
Corollary~\ref{cor:t}; \texttt{mem\_twice\_iff} and \texttt{DR\_iff}: the first equivalence in
Theorem~\ref{thm:struct}(b).
\item \texttt{corner} and \texttt{ten\_mem\_twice}: the corner of the array displayed in
\seqnum{A341260}, and $10\in\D$, as checks that the definitions match the entries.
\end{itemize}

The proof follows Sections~\ref{sec:circle} and~\ref{sec:gaps}, except that the full
descriptions of $\D$ and $\M$ by $T$ and $T'$ are not formalized; Lean uses only what the cells
need, namely the intervals $I_1,\dots,I_9$ explicitly and, for the remaining columns, bounds
such as $I_k\subset(0,\varphi^{-6})$ for even $k\ge8$ and $I_k\subset(-\varphi^{-9},0)$ for odd
$k\ge11$, which follow from $|\psi^k|\le|\psi|^K$ for $k\ge K$.
A program computes, for each cell and each $j$, the integer $c$ with
$t(m+j)=t(m)+j-c\varphi$ on the whole cell, and the lemma that decides membership of $m+j$; the
same data produce Tables~\ref{tab:twice} and~\ref{tab:missing}. Lean checks every resulting
inequality by linear arithmetic over the reals, using only the enclosure
$1.6180339<\varphi<1.618034$; the irrationality of $\varphi$ (in Mathlib) is used for the
boundary points excluded via \eqref{eq:irr}. An error in the program can therefore make a proof
fail but cannot make a false statement accepted. The proofs use no axioms beyond the three
standard ones of Lean's classical logic.

Not formalized: Theorems~\ref{thm:freq} and~\ref{thm:zeck}, and the full form of
Theorem~\ref{thm:struct}(b),(c). The array was also enumerated directly from its definition up to
$2\cdot10^9$, computing $\cl{n\varphi}$ exactly in integer arithmetic. The results agree with
Theorems~\ref{thm:main}, \ref{thm:freq} and~\ref{thm:zeck} on that range, and with the data in the
three OEIS entries.

\section{Remarks}\label{sec:remarks}

\begin{enumerate}
\item The entry \seqnum{A341260} notes that row $n$ of the ceiling array is a row of the extended
Wythoff array. We have not used this; Proposition~\ref{prop:col} plays the corresponding role.
\item By Theorem~\ref{thm:zeck}, membership in $\D$ and in $\M$ is a regular condition on
Zeckendorf representations, so both sets are Fibonacci-automatic. Statements such as
Theorem~\ref{thm:main}(b),(c) are then expressible in the first-order theory decided by Walnut
\cite{mousavi,shallitbook}, which Shallit \cite{shallit25} has used for related arrays. We have not
attempted such a proof; the argument here is independent of that approach.
\item The same analysis applies to the Wythoff array itself. There Lemma~\ref{lem:row} holds with
$\fl{n\varphi}$ in place of $\cl{n\varphi}$, so $u_n\in(\varphi-1,\varphi)$ and $I_k$ is replaced by
$\psi^k\cdot(\varphi-1,\varphi)$. One checks in the same way that the analogues of $E$ and $O$ are
disjoint and cover $(0,\varphi)$ apart from countably many points that are not of the form $t(m)$
with $m\ge1$, which recovers the fact that every positive integer occurs exactly once. The
ceiling shifts the window by one unit of $\varphi$, and this creates both the overlaps and the
holes.
\end{enumerate}

\section{Use of AI tools}\label{sec:ai}

This work was carried out with an AI system, Claude (Anthropic), under the direction of the
author. The AI system carried out the literature search, found the proofs, wrote the Lean
formalization and all computer programs, and wrote the text of this draft. The author directed
the work and chose among the options presented. The statements listed in
Section~\ref{sec:lean} as formalized have been checked by Lean; the remaining arguments have been
checked by exact computation as described there but have not yet been reviewed by a
mathematician. The literature search is recorded in the repository.

\begin{thebibliography}{99}

\bibitem{kimberling95}
C. Kimberling, The Zeckendorf array equals the Wythoff array, \emph{Fibonacci Quart.}
\textbf{33} (1995), 3--8.

\bibitem{kn}
L. Kuipers and H. Niederreiter, \emph{Uniform Distribution of Sequences}, Wiley, 1974.

\bibitem{link}
C. Link, Gaps in the intersection of two Beatty sequences: three conjectures of Kimberling,
draft, 2026, \url{https://github.com/clink-bw/math-proofs}.

\bibitem{mathlib}
The mathlib Community, The Lean mathematical library, in \emph{Proceedings of the 9th ACM
SIGPLAN International Conference on Certified Programs and Proofs}, ACM, 2020, pp.~367--381.

\bibitem{lean4}
L. de Moura and S. Ullrich, The Lean 4 theorem prover and programming language, in
\emph{Automated Deduction -- CADE 28}, Lecture Notes in Comput. Sci., Vol.~12699, Springer,
2021, pp.~625--635.

\bibitem{morrison}
D. R. Morrison, A Stolarsky array of Wythoff pairs, in \emph{A Collection of Manuscripts Related
to the Fibonacci Sequence}, Fibonacci Association, 1980, pp.~134--136.

\bibitem{mousavi}
H. Mousavi, Automatic theorem proving in Walnut, preprint, 2016,
\url{https://arxiv.org/abs/1603.06017}.

\bibitem{shallitbook}
J. Shallit, \emph{The Logical Approach to Automatic Sequences: Exploring Combinatorics on Words
with Walnut}, London Math. Soc. Lecture Note Ser., Vol.~482, Cambridge University Press, 2022.

\bibitem{shallit25}
J. Shallit, An ``experimental mathematics'' approach to Stolarsky interspersions via automata
theory, preprint, 2025, \url{https://arxiv.org/abs/2502.03312}.

\bibitem{oeis}
N. J. A. Sloane et al., The On-Line Encyclopedia of Integer Sequences, 2026,
\url{https://oeis.org}.

\bibitem{weyl}
H. Weyl, \"Uber die Gleichverteilung von Zahlen mod.\ Eins, \emph{Math. Ann.} \textbf{77}
(1916), 313--352.

\bibitem{zeckendorf}
E. Zeckendorf, Repr\'esentation des nombres naturels par une somme de nombres de Fibonacci ou de
nombres de Lucas, \emph{Bull. Soc. Roy. Sci. Li\`ege} \textbf{41} (1972), 179--182.

\end{thebibliography}

\bigskip
\noindent 2020 \emph{Mathematics Subject Classification}: Primary 11B39; Secondary 11B83, 11J71,
68V20.

\noindent \emph{Keywords}: Wythoff array, Fibonacci numbers, Zeckendorf representation, golden
ratio, rotation, formal proof.

\begin{sloppypar}
\noindent (Concerned with sequences \seqnum{A000045}, \seqnum{A001950}, \seqnum{A035513},
\seqnum{A341260}, \seqnum{A341261} and \seqnum{A341262}.)
\end{sloppypar}

\end{document}
